\begin{center}
{\Large\bfseries \en{UTI}}\\[4pt]
فصل ۱۴۰ هاریسون
\end{center}
\begingroup\small
این فایل رونویسی از ۹ صفحهٔ دست‌نویس است. شمارهٔ صفحه و نام تصویر، تطبیق با اصل را ممکن می‌کند. عبارت‌های قرمز، خوانش نامطمئن دست‌خط هستند؛ هایلایت‌های زرد از اصل یادداشت حفظ شده‌اند.
\endgroup
\section*{صفحهٔ ۱}
{\footnotesize\en{20261005\_001146641.jpg}}

شایع‌ترین تظاهر بالینی \en{UTI}، سیستیت حاد است.

وجود باکتری در مجاری ادراری بدون وجود علائم مرتبط، مطرح‌کنندهٔ \en{UTI} نیست.
\begin{itemize}
\item \en{ASB} (باکتری ادراری بدون علامت).
\item \en{ASB} در اغلب موارد هیچ خطری برای فرد ندارد و نیاز به درمان وجود ندارد.
\item تنها دو گروه از افراد با \en{ASB} را باید درمان کنیم:
\begin{itemize}
\item زنان باردار.
\item افرادی که می‌خواهند تحت پروسیجرهای اورولوژیک قرار گیرند؛ پروسیجرهایی که به مخاط آسیب می‌زنند.
\end{itemize}
\item عوارض \en{ASB} در زنان باردار: پیلونفریت، زایمان پره‌ترم، \en{LBW}.
\end{itemize}

\begin{itemize}
\item \en{Cystitis}: عفونت علامت‌دار مثانه است.
\item \en{Prostitis}: عفونت علامت‌دار پروستات است. % Spelling as handwritten on page 1.
\item \en{Pyelonephritis}: عفونت علامت‌دار کلیه.
\item \en{Uncomplicated UTI}: عفونت محدود به مثانه در زنان و مردان بدون کاتتر ادراری.
\item \en{Recurrent UTI} لزوماً به معنای \en{Complicated} بودن نیست.
\end{itemize}

پزشک در قدم اول باید تشخیص دهد:
\begin{itemize}
\item بیمار پایدار است و می‌تواند سرپایی درمان شود یا باید بستری گردد.
\item آیا نشانه‌ای برای عفونت راجعه وجود دارد؟ سنگ، انسداد.
\item آیا \en{AB} تجویز‌شده در فضای کلیه و ادرار به غلظت مناسب می‌رسد یا خیر.
\end{itemize}

پروستات:
\begin{itemize}
\item باکتریال حاد (\en{ABP}).
\item باکتریال مزمن (\en{CBP}).
\item منشأ: عارضهٔ سیستیت باشد.
\item تظاهر: هماتوژن استاف اورئوس باشد در مردان با \unclear{\en{And few}}.
% The final Latin phrase in this line is not confidently legible.
\end{itemize}

گایدلاین‌های بالینی درمان \en{UTI} عارضه‌دار قابل تعمیم برای درمان پروستاتیت نیستند.

\sourcepage{۲}{20261005\_001156435.jpg}
پروستاتیت اغلب در زمینهٔ نایسریا گنوره، کلامیدیا تراکوماتیس، مایکوپلاسما ژنیتالیوم رخ می‌دهد ولی می‌تواند در زمینهٔ \en{E. Coli} نیز باشد.

\subsection*{اپیدمیولوژی}
\begin{itemize}
\item در نوزادی شیوع \en{UTI} اندکی در پسرها بیشتر است.
\item پس از ۵۰ سالگی، به علت \en{BPH}، شیوع \en{UTI} در مردان و زنان برابر است.
\item اغلب \en{UTI}ها در مردان به علت سونداژ یا آنومالی‌ها می‌باشد.
\item شیوع \en{ASB} در زنان \rng{20}{40} سال در حدود ۵\% است.
\item \rng{40}{50}\% در مردان و زنان مسن.
\item تا ۸۰\% زنان حداقل یک نوبت \en{UTI} را در زندگی تجربه می‌کنند؛ اغلب سیستیت ساده.
\item شیوع حملهٔ \en{UTI} در مردان ۱۴\% است.
\item \rng{20}{30}\% زنانی که یک نوبت \en{UTI} داشته‌اند دچار \en{UTI} راجعه می‌شوند.
\item \en{Early Recurrence}: ظرف ۲ هفته؛ عود در نظر گرفته می‌شود. نیازمند بررسی علت است.
\item چندین اپیزود \en{UTI}: \en{Clustering}؛ ممکن است یک \en{RF} جدید به وجود آمده باشد.
\item سن زیر ۱۵ سال برای اولین \en{UTI}، یک \en{RF} برای \en{UTI} راجعه است.
\item برخی از افراد استعداد ابتلا به \en{UTI} بیشتری دارند.
\end{itemize}

\subsection*{ریسک فاکتورهای \en{UTI}}
\begin{itemize}
\item عواملی که جریان ادرار طبیعی را مختل می‌کنند.
\item عواملی که باکتری‌های جدید به فلور قبلی وارد می‌کنند.
\item عواملی که دفاع مخاطی طبیعی را به صورت موضعی تخریب می‌کنند.
\end{itemize}

\subsection*{اتیولوژی \en{UTI}}
\begin{itemize}
\item \en{E. Coli} شایع‌ترین پاتوژن در همهٔ تقسیم‌بندی‌ها است.
\item \en{Acute Cystitis}: \rng{75}{90}\% \en{E. Coli}؛ \rng{5}{15}\% استاف ساپروفیتیکوس؛ سایر انتروباکترها: کلبسیلا، پروتئوس.
\item پاتوژن‌های پیلونفریت مشابه سیستیت هستند؛ هرچند پیلونفریت اغلب ناشی از صعود سیستیت است.
\item \en{CAUTI}: \en{E. Coli} شایع‌ترین است؛ سودومونا آئروژینوزا، انتروکوک، استاف اورئوس، کاندیدا.
\end{itemize}

\sourcepage{۳}{20261005\_001203158.jpg}
مقاومت آنتی‌بیوتیکی در گونه‌های انتروباکتر در حال افزایش است.

\subsection*{رویکرد بالینی}
مهم‌ترین سؤال‌هایی که پس از شک به \en{UTI} باید پاسخ بدهیم:
\begin{itemize}
\item آیا بیمار واقعاً مبتلا به \en{UTI} می‌باشد؟
\item اگر مبتلا است آیا بیماری محدود به مثانه است یا نه؟
\item اگر تشخیص قطعی نیست آیا شروع آزمایشی درمان \en{AB} امپریکال ایمن است؟
\end{itemize}

\begin{itemize}
\item \en{ASB}: بیمار کشت ادراری مثبت دارد ولی علائم و نشانه‌های لوکالیزه‌کنندهٔ ادراری که با چیز دیگری توجیه شوند ندارد.
\item خطای تشخیصی باعث می‌شود \en{ASB} پیدا کنیم که بعد در افتراق آن از \en{UTI} دچار چالش شویم.
\item بجز در موارد غربالگری، در بیمار بدون علامت \en{U/C} درخواست نمی‌کنیم.
\item وجود پیوری در فرد بدون علامت نقشی در تشخیص \en{ASB} از \en{UTI} ایجاد نمی‌کند.
\item پیوری بسیار شدید می‌تواند در استاز ادراری نیز دیده شود؛ همچنین در \en{ESRD}.
\end{itemize}

\subsection*{\en{Cystitis}}
\begin{itemize}
\item علائم تیپیک: دیزوری، فرکوئنسی، اورجنسی.
\item سایر علائم: سوپراپوبیک دیس‌کامفورت یا تندرنس؛ ناکچوری اخیر، \en{hesitancy}، هماچوری گروس.
\item علائم نامنطبق: تب، \en{rigors}، درد فلانک؛ مطرح‌کنندهٔ گسترش عفونت.
\item در زنان می‌توان صرفاً بر مبنای شرح حال درمان کرد.
\item در زنی با علائم سیستیت که \en{U/A} منفی است، بیماری \en{R/O} نمی‌شود؛ \en{U/C}.
\item مواردی که در زنان حتماً به \en{U/C} نیاز داریم: زنان باردار، زنان مشکوک به میکروارگانیسم‌های مقاوم، زنان با \en{UTI} راجعه.
\item علائم و نشانه‌های سیستیت در مردان مشابه زنان است.
\item در همهٔ موارد \en{UTI} در مردان باید کشت ادرار (\en{U/C}) انجام شود.
\end{itemize}

\subsection*{\en{Prostatitis}}
\begin{itemize}
\item شامل اتیولوژی‌های عفونی و غیرعفونی غدهٔ پروستات می‌شود.
\item عفونت می‌تواند حاد یا مزمن باشد.
\item تقریباً همواره باکتریال است: در فرم عفونی.
\item فرم غیرعفونی بسیار شایع‌تر است: \en{Chronic Pelvic Pain Syndrome (CPPS)}.
\end{itemize}

\sourcepage{۴}{20261005\_001210587.jpg}
\begin{itemize}
\item \en{ABP}: دیزوری، فرکوئنسی، درد در ناحیهٔ آنال.
\item می‌تواند با تب، \en{rigors} و انسداد ادراری همراه باشد.
\item \en{CBP}: روند تدریجی دارد؛ به صورت اپیزودهای مکرر سیستیت شروع می‌شود.
\end{itemize}
\begin{itemize}
\item علائم تیپیک عفونت معمولاً وجود ندارند.
\item در برخی موارد درد وجود دارد.
\item مردان با سیستیت راجعه باید از نظر پروستاتیت و ریتنشن ادراری بررسی شوند.
\item \en{UTI} تب‌دار در مردان می‌تواند باعث افزایش \en{PSA} شود و پروستات و غدد سمینال وزیکول در سونوگرافی افزایش سایز خواهند داشت.
\item مردان با \en{UTI} تب‌دار حتماً باید تصویربرداری انجام دهند: سونوگرافی یا \en{CT}.
\item در مردان اگر تشخیص \en{UTI} غیرقطعی است یا \en{UTI} راجعه می‌باشد: ارجاع به اورولوژیست.
\end{itemize}

\subsection*{پیلونفریت}
\begin{itemize}
\item تب عامل افتراقی بین سیستیت و پیلونفریت است.
\item پیلونفریت خفیف می‌تواند تب \en{low-grade} با یا بدون درد فلانک باشد.
\item در بیماران انسدادی می‌تواند منجر به \en{acute papillary necrosis} شود که می‌تواند پیلونفریت را عارضه‌دار کند.
\item پیلونفریت آمفیزماتو: تقریباً همیشه در بیماران با \en{DM} با کنترل ضعیف رخ می‌دهد.
\item پیلونفریت زانتوگرانولوماتوز: انسداد مزمن در کنار عفونت مزمن منجر به عفونت چرکی می‌شود که بافت کلیه را تخریب می‌کند.
\item تشکیل آبسهٔ اینترا یا پرینفرال می‌تواند پیلونفریت را عارضه‌دار کند.
\item تب ادامه‌دار $\pm$ باکتری ادراری، علی‌رغم \en{ABT}.
\end{itemize}

\subsection*{\en{CAUTI}}
\begin{itemize}
\item علائم + باکتری ادراری + پیوری + کاتتر.
\item حداقل تعداد باکتری باید $\geq 10^3$ کلونی باشد.
\item باکتری ادراری به دنبال کاتتر غیرقابل اجتناب است و مشکل تشخیصی همین‌جاست.
\item ارزش تشخیصی همهٔ علائم \en{UTI} برای \en{CAUTI} کمتر است؛ حتی لکوسیتوز.
\item حتماً باید سایر علل علائم موجود \en{R/O} شود.
\end{itemize}

\sourcepage{۵}{20261005\_001219869.jpg}
\subsection*{ابزارهای تشخیصی}
\begin{itemize}
\item خودتشخیصی در زنان مبتلا به \en{UTI} راجعه تا حدود زیادی دقیق است.
\item در زنان با حداقل یک علامت لوکالیزه‌کننده (دیزوری و ...) و بدون عامل عارضه‌دارکنندهٔ دیگر، احتمال سیستیت حاد یا پیلونفریت ۵۰\% است.
\item همین بیمار اگر ترشح واژینال نداشته باشد، \en{RF} برای \en{UTI} داشته باشد، احتمال ۹۰\% می‌رود.
\item \en{U/A} و \en{U/C} دیگر ضروری نیستند، مگر شک به مقاومت.
\item خطر این روش این است که بیمار کلامیدیا تراکوماتیس به جای \en{UTI} درمان می‌شود.
\item وجود ترشح واژینال قویاً احتمال \en{STI} را مطرح می‌کند.
\item می‌توان قبل از درمان به عنوان \en{UTI}، غربالگری \en{STI} انجام داد.
\end{itemize}

\subsection*{\en{U/A}، \en{U/C}، \en{Dipstick}}
تفاوت:
\begin{itemize}
\item \en{D/S} برای بررسی \en{LE} و \en{N}؛ \en{U/A} برای تعداد \en{RBC} و \en{WBC}؛ \en{U/C} اثبات وجود باکتری و الگوی مقاومت.
\item همهٔ بیماران سونداژ هم پیوری دارند هم باکتری ادراری؛ اگر علامت نداشته باشند، به معنی \en{UTI} نیست.
\item فقط باکتری‌های خانوادهٔ انتروباکتریاسه نیترات را به نیتریت تبدیل می‌کنند.
\item باید مقدار کافی نیتریت در ادرار باشد تا بتوان آن را شناسایی نمود.
\item نیتریت در بیمار با \en{ASB} انتروباکتر مثبت است.
\item اگر بیمار فرکوئنسی داشته باشد مقدار کافی ادرار جمع نمی‌شود که نیتریت مثبت گردد.
\item \en{D/S} منفی برای \en{LE} و \en{N} عملاً به معنی رد \en{UTI} است.
\item \en{D/S} منفی نباید برای رد وجود باکتری ادراری در زن باردار ملاک قضاوت باشد.
\end{itemize}

\subsection*{میکروسکوپی ادرار}
\begin{itemize}
\item در \rng{90}{100}\% بیماران سیستیت، پیوری را می‌بینیم و در ۳۰\% آن‌ها هماچوری.
\item پیوری یعنی $\mathrm{WBC}>10$ در هر \en{HPF}؛ البته در زنان مسن باید خیلی بالاتر باشد: $>250$.
\end{itemize}

\subsection*{\en{U/C}}
\begin{itemize}
\item هم در \en{UTI} و هم در \en{ABB} مثبت می‌شود. % ABB is the handwritten abbreviation in this line.
\item در زنان با علائم سیستیت تعداد کلونی باید $>10^2$ باشد (حساسیت ۹۰\%).
\item اگر \en{mixed bacterial} باشد مطرح‌کنندهٔ آلودگی است، مگر آنکه \en{long-term cath} باشد.
\end{itemize}
\textit{[خط‌خورده: درمان \en{UTI}]} % Crossed-out heading on source page 5.

\sourcepage{۶}{20261005\_001228230.jpg}
اپروچ تشخیصی برای \en{UTI} / تظاهر بالینی / ویژگی‌های بیمار / تشخیص و درمان

\subsection*{۱) شروع ناگهانی علائم ادراری (دیزوری، فرکوئنسی، اورجنسی)}
\begin{itemize}
\item زن سالم غیر باردار با ریسک پایین برای مقاومت آنتی‌بیوتیکی:
\begin{itemize}
\item درمان به عنوان سیستیت بدون عارضه.
\item نیاز به \en{U/C} ندارد.
\end{itemize}
\item زن با شرح حال یا \en{RF} برای \en{STI}:
\begin{itemize}
\item یا سیستیت بدون عارضه یا \en{STI}.
\item \en{U/A}، \en{U/C}، ارزیابی \en{STI}، معاینهٔ لگن.
\end{itemize}
\item مرد با درد پرینه، لگن یا پروستات:
\begin{itemize}
\item پروستاتیت حاد مطرح است.
\item \en{U/A}، \en{U/C}، ارجاع به اورولوژیست.
\end{itemize}
\item بیمار با کاتتر ادراری:
\begin{itemize}
\item \en{CAUTI} مطرح است.
\item تعویض یا برداشتن کاتتر.
\item \en{U/A}، \en{U/C}.
\item اگر تب دارد، \en{B/C}.
\end{itemize}
\item سایر بیماران:
\begin{itemize}
\item احتمال \en{UTI} عارضه‌دار.
\item \en{U/A}، \en{U/C}، ارزیابی آبنورمالیتی‌های آناتومیک یا عملکردی.
\end{itemize}
\end{itemize}

\subsection*{۲) شروع ناگهانی درد، تهوع و استفراغ با تب $\pm$ علائم سیستیت}
\begin{itemize}
\item زن سالم غیر باردار:
\begin{itemize}
\item پیلونفریت مطرح است.
\item \en{U/A}، \en{U/C}، درمان سرپایی.
\end{itemize}
\item سایر بیماران:
\begin{itemize}
\item پیلونفریت یا پروستاتیت حاد (در مردان).
\item \en{U/C}، \en{B/C}.
\end{itemize}
\end{itemize}

\subsection*{۳) علائم سیستمیک (تب، \en{WBC}، لکوسیتوز)}
\begin{itemize}
\item بیماران مسن، آسیب نخاعی، نقص ایمنی؛ هیچ تشخیص دیگری وجود ندارد.
\item \en{UTI} عارضه‌دار مطرح است.
\item \en{U/C}، \en{B/C}، بررسی سایر علل احتمالی.
\end{itemize}

\sourcepage{۷}{20261005\_001234198.jpg}
\subsection*{۴) بدون علائم ادراری}
\begin{itemize}
\item \en{U/C} مثبت در بیمار باردار یا بیماری که قصد دارد پروسیجر اورولوژیک انجام دهد:
\begin{itemize}
\item \en{ASB} مطرح است.
\item درمان انجام شود.
\end{itemize}
\item \en{U/C} مثبت در بقیهٔ موارد:
\begin{itemize}
\item \en{ASB} مطرح است.
\item درمان لازم نیست.
\end{itemize}
\item \en{U/C} مثبت در فرد با کاتتر ادراری:
\begin{itemize}
\item \en{CA-ASB} مطرح است.
\item نیازمند بررسی و درمان نیست.
\item کاتترهای اضافی حذف شوند.
\end{itemize}
\end{itemize}

\subsection*{۵) علائم ادراری حاد راجعه}
\begin{itemize}
\item زن سالم غیر باردار:
\begin{itemize}
\item سیستیت راجعه مطرح است.
\item \en{U/C} برای تأیید تشخیص.
\item پروفیلاکسی یا \en{patient initiated management}.
\end{itemize}
\item مردان:
\begin{itemize}
\item \en{CBP} مطرح است.
\item ارجاع به اورولوژیست.
\end{itemize}
\end{itemize}

\subsection*{درمان عفونت مجاری ادراری}
\begin{itemize}
\item \hl{همهٔ \en{UTI} حقیقی نیازمند درمان \en{ABT} است.}
\item سیستیت بدون عارضه در مردان باید \hl{۷ روزه} \en{ABT} شود.
\item \en{TMP-SMX}، نیتروفورانتوئین، فسفومایسین، \en{pivmecillinam}.
\item برای \en{E. Coli MDR}: فسفومایسین، پیومسیلینام.
\item خط‌های بعدی: فلوروکینولون، $\beta$-لاکتام (\rng{5}{7} روز).
\item در سیستیت بدون عارضه در زنان فلوروکینولون نمی‌دهیم، مگر آنکه گزینهٔ دیگری نداشته باشیم.
\item در مردان اگر شک به پروستاتیت وجود دارد، فلوروکینولون و \en{TMP-SMX} بر $\beta$-لاکتام‌ها ارجح هستند چون نفوذ بهتری دارند.
\item درمان سیستیت بدون عارضه در زنان: دورهٔ ۵ روزهٔ نیتروفورانتوئین و دورهٔ ۳ روزهٔ \en{TMP-SMX} به یک اندازه مؤثر هستند.
\item نیتروفورانتوئین در $\mathrm{GFR}<30$ ممنوع است.
\item در درمان پروستاتیت و پیلونفریت جایگاهی ندارد.
\item به دلیل ریسک سمیت ریوی در پروفیلاکسی سیستیت راجعه جایگاهی ندارند.
\end{itemize}

\sourcepage{۸}{20261005\_001244694.jpg}
\begin{itemize}
\item فسفومایسین علیه \en{E. Coli} عالی است ولی علیه انتروباکترها قطعی نیست.
\item در درمان پیلونفریت جایگاهی ندارد.
\item رفع \en{discomfort} در سیستیت: \hl{فنازوپیریدین} (فقط در مجاری ادراری)؛ عارضهٔ تهوع دارد.
\end{itemize}

\subsection*{استراتژی‌های درمانی در سیستیت حاد بدون عارضه}
% F, M, BD, DS, and alpha symbols preserved from handwritten rows.
\begingroup
\small
\setlength{\tabcolsep}{4pt}
\renewcommand{\arraystretch}{1.35}
\begin{longtable}{@{}p{.20\linewidth}p{.20\linewidth}p{.25\linewidth}p{.27\linewidth}@{}}
\toprule
دارو & دوز & مدت & عوارض\\
\midrule
\endfirsthead
\toprule
دارو & دوز & مدت & عوارض\\
\midrule
\endhead
نیتروفورانتوئین & \en{100 mg BD} & \en{M, F}: \hl{\rng{5}{7} روز} & تهوع، سردرد\\
\en{TMP-SMX} & \en{1 DS tab BD} & \en{F}: ۳ روز؛ \en{M}: ۷ روز & راش، آنمی، تهوع، هایپرکالمی\\
فسفومایسین & \en{3 gr Sachet} & \en{F}: ۱ روز؛ \en{M}: ۳ روز & اسهال، تهوع، سردرد\\
پیومسیلینام & \en{400 mg BD} & \en{F}: \rng{3}{7} روز & تهوع، استفراغ، اسهال\\
فلوروکینولون & $\alpha$ & \en{F}: ۳ روز؛ \en{M}: ۷ روز & تهوع، بی‌خوابی، سردرد\\
بتالاکتام & $\alpha$ & \en{M, F}: \rng{5}{7} روز & اسهال، راش، کاند\\
\bottomrule
\end{longtable}
\endgroup

\subsection*{درمان پروستاتیت}
\begin{itemize}
\item در شک به \en{ABP}، پس از اخذ \en{U/C} و \en{B/C} بلافاصله \en{ABT} را شروع می‌کنیم.
\item درمان \en{ABP}: \hl{حداقل ۲ هفته} و \hl{ترجیحاً ۴ هفته}، فلوروکینولون یا \en{TMP-SMX}.
\item تشخیص آبسهٔ پروستات: ارجاع به اورولوژیست.
\item درمان \en{CBP}: \hl{\rng{4}{6} هفته} درمان \en{ABT} و در صورت راجعه‌شدن، \hl{۱۲ هفته} \en{ABT}.
\end{itemize}

\subsection*{درمان پیلونفریت}
\begin{itemize}
\item درمان امپریکال انتخابی \hl{فلوروکینولون}ها هستند؛ \hl{\rng{5}{7} روز}.
\item درمان در بیمار بستری بتالاکتام وریدی است: سفتریاکسون، سفپیم.
\item در بیمار سرپایی سیپرو، \en{TMP-SMX}، لووفلوکساسین، تازوسین، کارباپنم‌ها؛ \rng{7}{14} روز.
\item خط‌های جایگزین: آمینوگلیکوزیدها.
\item اگر \hl{تب ظرف ۷۲ ساعت} از شروع \en{ABT} قطع نشد: بررسی.
\item \en{XGP} اغلب به نفرکتومی می‌انجامد.
\end{itemize}

\subsection*{درمان \en{CAUTI}}
\begin{itemize}
\item \hl{\rng{7}{14} روز} بر اساس آنتی‌بیوگرام.
\item کاتتر باید تعویض شود.
\end{itemize}

\sourcepage{۹}{20261005\_001251895.jpg}
\subsection*{درمان \en{UTI} در زنان باردار}
\begin{itemize}
\item آمینوپنی‌سیلین‌ها (آمپی‌سیلین، آموکسی‌سیلین)، سفالوسپورین‌ها استفاده می‌شوند.
\item اگر هر بتالاکتام نتوانیم بدهیم، نیتروفورانتوئین می‌دهیم؛ در تریمستر اول کمی خطرناک.
\item \hl{سولفونامیدها (\en{TMP-SMX})} در سه‌ماههٔ اول \hl{ممنوع هستند}.
\item در \en{near term} نیز \hl{ممنوع هستند}: خطر کرن‌ایکتروس.
\item \hl{فلوروکینولون‌ها} در بارداری \hl{ممنوع هستند}: اختلال در ساخت غضروف در جنین.
\item طول مدت درمان زنان باردار \hl{\rng{5}{7} روز} است.
\item در پیلونفریت در بارداری: بتالاکتام وریدی $\pm$ آمینوگلیکوزید.
\end{itemize}

\subsection*{کاندیدوری}
\begin{itemize}
\item \en{RF}: کاتتر ادراری، بستری در \en{ICU}، آنتی‌بیوتیک وسیع‌الطیف، دیابت زمینه‌ای.
\item می‌تواند بدون علامت تا سپسیس باشد.
\item برداشتن کاتتر ادراری در $1/3$ بیماران بدون علامت کاندیدوری را رفع می‌کند.
\item موارد بدون علامت را درمان نمی‌کنیم.
\item \hl{فلوکونازول \rng{200}{400} میلی‌گرم روزانه برای \rng{7}{14} روز}.
\item خط بعدی: فلوسیتوزین، آمفوتریسین \en{B}.
\end{itemize}

\subsection*{جلوگیری از بازگشت \en{UTI} در زنان}
\begin{itemize}
\item معیار ۳ نوبت و بیشتر \hl{\en{UTI} حقیقی در سال} می‌باشد ولی برای هر بیمار جداگانه تصمیم می‌گیریم.
\item تأیید تشخیص \en{UTI} فعلی، تصویربرداری برای پیلونفریت، سنگ، ارزیابی ریتنشن ادراری.
\item افزایش مصرف مایعات در افرادی که منع ندارند.
\item محصولات \en{Cranberry} مؤثر هستند در پیشگیری از بازگشت \en{UTI}.
\item درمان \en{ABT} پروفیلاکسی:
\begin{itemize}
\item برای ۶ ماه دارو می‌دهیم؛ قطع دارو؛ بازگشت؛ دوباره ۶ ماه: \en{low dose TMP-SMX}، نیتروفورانتوئین.
\item \en{ABT} به صورت \en{postcoital}.
\end{itemize}
\end{itemize}
